\section{Erd\H{o}s Problem \#200: longest arithmetic progression of primes up to $N$ (ROUND 2)}

\subsection*{1) ROUND-2 OBJECTIVE}
I again pursue path \textbf{(C)}.  This is Erd\H{o}s Problem \#200 \cite{Bloom200}.  Round 1 proved the standard upper bound
\[
L(N)\le (1+o(1))\log N,
\]
but did not isolate where that argument fundamentally stops.  In this round I sharpen the modular argument into an exact dichotomy for the first term of a prime progression, prove a stronger upper bound in the generic case, and then give an explicit family of \emph{pseudoprogressions} showing that the primorial/congruence strategy cannot by itself prove $L(N)=o(\log N)$.

\subsection*{2) Round-1 FOUNDATION USED}
From Round 1 I use the already-proved fact that if
\[
a,a+d,\dots,a+(k-1)d
\]
is a $k$-term arithmetic progression of primes, then every prime $q\le k/2$ divides $d$.  Equivalently,
\[
P(k/2):=\prod_{q\le k/2} q \mid d.
\]
I also use the Round-1 bound $k\le (1+o(1))\log N$.

\subsection*{3) NEW INSIGHT / TOOL (ROUND 2)}
The new tool is a precise divisibility dichotomy.

\medskip
\noindent\textbf{Key new dichotomy.}
For all sufficiently large $k$, either
\begin{itemize}
\item the first prime $a$ satisfies $a>k$, in which case \emph{every} prime $q\le k$ divides $d$; or
\item the first prime $a$ itself is the only possible exception, and then necessarily
\[
a\in (k/2,k]\cap \mathbf P
\]
and
\[
P(k)/a\mid d.
\]
\end{itemize}
So the only obstruction to improving the Round-1 logarithmic upper bound by the primorial method is the exceptional regime where the first term itself is a prime in $(k/2,k]$.

\subsection*{4) ATTACK PLAN (ROUND 2)}
The exact Round-1 gap was that the modular argument only yielded $d\gg e^{(1+o(1))k}$, which still allows $k\asymp \log N$.

I now proceed in two steps.
\begin{enumerate}
\item Strengthen the divisibility statement from primes $\le k/2$ to primes $\le k$, up to one forced exceptional prime.
\item Show that if the exception does not occur (that is, if $a>k$), then the extra factor of $k$ in $(k-1)P(k)$ gives the stronger upper bound
\[
k\le \log N-\log\log N+o(\log\log N).
\]
Then I explicitly construct pseudoprogressions at scale $e^{(1+o(1))k}$ that satisfy all local congruence conditions up to $k$, proving that those local conditions alone can never yield $o(\log N)$.
\end{enumerate}

\subsection*{5) WORK (ROUND 2)}
Let
\[
p_j=a+jd\qquad (0\le j\le k-1)
\]
be a $k$-term arithmetic progression of primes, with $p_{k-1}\le N$.
Write
\[
P(x):=\prod_{q\le x} q,
\qquad
\vartheta(x):=\sum_{q\le x}\log q,
\]
where in both places $q$ runs over primes.

\medskip
\noindent\textbf{Proposition 5.1.}
If $a\le k$, then in fact $a\in (k/2,k]$.

\smallskip
\noindent\emph{Proof.}
The first term $a$ is prime.  If $a\le k/2$, then by the Round-1 divisibility lemma the prime $a$ divides $d$.  Hence
\[
a+d
\]
is divisible by $a$ and is strictly larger than $a$, so it is composite.  This contradicts the primality of the second term.  Therefore $a>k/2$. \hfill$\square$

\medskip
\noindent\textbf{Theorem 5.2 (exact divisibility dichotomy).}
For all sufficiently large $k$ (equivalently, whenever $P(k/2)>k$), exactly one of the following two alternatives holds:
\begin{enumerate}
\item $a>k$ and $P(k)\mid d$;
\item $a\in (k/2,k]$ is prime and $P(k)/a\mid d$.
\end{enumerate}

\smallskip
\noindent\emph{Proof.}
By the Round-1 lemma,
\[
P(k/2)\mid d.
\]
For sufficiently large $k$ this implies $d\ge P(k/2)>k$.

Now let $q$ be any prime with $k/2<q\le k$.  Suppose $q\nmid d$.  Since $d$ is invertible modulo $q$, the congruence
\[
a+jd\equiv 0\pmod q
\]
has exactly one residue class of solutions modulo $q$, so among the $k$ indices $0,1,\dots,k-1$ there is at least one $j$ with $a+jd\equiv 0\pmod q$.  Because every term of the progression is prime, that term must equal $q$.

But if $j\ge 1$, then
\[
a+jd\ge a+d>d>k\ge q,
\]
so $a+jd$ cannot equal $q$.  Therefore the only possible divisible term is the first one $j=0$, and hence $q=a$.

Thus every prime $q\le k$ divides $d$, except possibly the single prime $q=a$.
If $a>k$, there is no exception at all, so $P(k)\mid d$.
If $a\le k$, then Proposition 5.1 shows that $a\in (k/2,k]$, and the only possible missing factor is exactly $a$, so $P(k)/a\mid d$.
This proves the dichotomy. \hfill$\square$

\medskip
\noindent\textbf{Corollary 5.3 (generic-case strengthening).}
Assume $a>k$.  Then
\[
k\le \log N-\log\log N+o(\log\log N).
\]

\smallskip
\noindent\emph{Proof.}
In the case $a>k$, Theorem 5.2 gives $P(k)\mid d$, so $d\ge P(k)$.  Therefore
\[
N\ge a+(k-1)d>(k-1)P(k).
\]
Taking logarithms,
\[
\log N\ge \log(k-1)+\vartheta(k).
\]
By the classical de la Vall\'ee Poussin form of the prime number theorem,
\[
\vartheta(k)=k+O\bigl(k e^{-c\sqrt{\log k}}\bigr)
\]
for some absolute $c>0$.  The error term is $o(\log k)$, so
\[
\log N\ge k+\log k+o(\log k).
\]
Round 1 already gave $k\le (1+o(1))\log N$, hence
\[
\log k=\log\log N+O(1).
\]
Substituting this into the previous display gives
\[
k\le \log N-\log\log N+o(\log\log N),
\]
as claimed. \hfill$\square$

\medskip
\noindent\textbf{Proposition 5.4 (explicit pseudoprogressions).}
For every integer $k\ge 2$ there exists a length-$k$ arithmetic progression
\[
x_j=p+jP(k)\qquad (0\le j\le k-1)
\]
with the following properties:
\begin{enumerate}
\item $p$ is prime and $k<p\le 2k$;
\item no term $x_j$ is divisible by any prime $q\le k$;
\item the last term satisfies
\[
x_{k-1}\le 2k+(k-1)P(k),
\]
so in particular
\[
\log x_{k-1}=k+O(\log k).
\]
\end{enumerate}

\smallskip
\noindent\emph{Proof.}
By Bertrand's postulate there exists a prime $p$ with $k<p\le 2k$.
Set
\[
d:=P(k).
\]
For any prime $q\le k$, we have $q\mid d$, hence
\[
x_j=p+j d\equiv p\not\equiv 0\pmod q,
\]
because $p>k\ge q$.  Thus no $x_j$ has a prime divisor at most $k$.
The bound on $x_{k-1}$ is immediate from the definition, and
\[
\log P(k)=\vartheta(k)=k+o(k)
\]
by the prime number theorem, which implies
\[
\log x_{k-1}=k+O(\log k).
\]
\hfill$\square$

\medskip
\noindent\textbf{Consequence for proof strategies.}
Proposition 5.4 shows that local congruence information modulo primes $q\le k$ allows length-$k$ arithmetic progressions all the way up to scale $e^{(1+o(1))k}$.  Therefore:
\begin{quote}
\emph{Any argument based only on congruence obstructions modulo primes up to $k$ cannot prove $L(N)=o(\log N)$.}
\end{quote}
Such an argument can force at most the scale $k\ll \log N$, and Round 1 already reaches that barrier.

\subsection*{6) ADVERSARIAL VERIFICATION}
I now test the new work for hidden weaknesses.
\begin{itemize}
\item \emph{Is the exceptional case real, or just an artefact of the proof?}  It is real.  The prime progression
\[
5,11,17,23,29
\]
has $k=5$, first term $a=5\in (k/2,k]$, and common difference $d=6$, which is not divisible by $5$.
\item \emph{Did Proposition 5.4 accidentally prove something about prime progressions?}  No.  It proves only the existence of arithmetic progressions avoiding \emph{small} prime divisors.  That is exactly the intended barrier: local congruence conditions do not distinguish these pseudoprogressions from genuinely prime progressions.
\item \emph{Does Theorem 5.2 require large $k$?}  Yes, because I used $P(k/2)>k$ to deduce $d>k$.  This is harmless for an asymptotic problem; finitely many small values can be absorbed into constants.
\item \emph{Did Corollary 5.3 use more than was proved?}  No.  It uses only $N>(k-1)P(k)$ and the standard error term for $\vartheta(k)$.
\item \emph{Does the round really sharpen Round 1?}  Yes.  It isolates the only local exceptional regime, proves a stronger upper bound when that regime is absent, and rigorously explains why the old primorial method stops where it stops.
\end{itemize}

\subsection*{7) FINAL}
\textbf{UNRESOLVED (BUT STRICTLY ADVANCED).}

Round 2 proves two genuinely new things beyond Round 1.
\begin{enumerate}
\item For long prime progressions, the only possible failure of full divisibility $P(k)\mid d$ is that the first term itself is a prime in $(k/2,k]$.
\item Outside that exceptional case one has the stronger bound
\[
k\le \log N-\log\log N+o(\log\log N).
\]
Moreover, explicit pseudoprogressions show that the purely local congruence/primorial strategy cannot by itself do better than the Round-1 logarithmic bound.
\end{enumerate}
So the remaining difficulty has been sharply localized: any proof of $L(N)=o(\log N)$ must either exploit information beyond local congruence obstructions, or else must somehow rule out the exceptional regime $a\in (k/2,k]$ for genuinely long prime progressions.

\subsection*{8) COMPLETION ESTIMATE (MANDATORY)}
\textbf{COMPLETION: 57\%}

\bigskip
\section*{References}

\begin{thebibliography}{99}

\bibitem{Bloom152}
T. F. Bloom,
\emph{Erd\H{o}s Problem \#152},
\texttt{https://www.erdosproblems.com/152}, accessed 2026-03-07.

\bibitem{Bloom200}
T. F. Bloom,
\emph{Erd\H{o}s Problem \#200},
\texttt{https://www.erdosproblems.com/200}, accessed 2026-03-07.

\bibitem{Bloom201}
T. F. Bloom,
\emph{Erd\H{o}s Problem \#201},
\texttt{https://www.erdosproblems.com/201}, accessed 2026-03-07.

\bibitem{ESS94}
P. Erd\H{o}s, A. S\'ark\"ozy, and T. S\'os,
\emph{On sum sets of Sidon sets. I},
Journal of Number Theory \textbf{47} (1994), 329--347.

\bibitem{GT08}
B. Green and T. Tao,
\emph{The primes contain arbitrarily long arithmetic progressions},
Annals of Mathematics \textbf{167} (2008), 481--547.

\bibitem{KSS75}
J. Koml\'os, M. Sulyok, and E. Szemer\'edi,
\emph{Linear problems in combinatorial number theory},
Acta Mathematica Academiae Scientiarum Hungaricae \textbf{26} (1975), 113--121.

\bibitem{MV07}
H. Montgomery and R. Vaughan,
\emph{Multiplicative Number Theory I. Classical Theory},
Cambridge Studies in Advanced Mathematics 97, Cambridge University Press, 2007.

\bibitem{OBryant04}
K. O'Bryant,
\emph{A complete annotated bibliography of work related to Sidon sequences},
Electronic Journal of Combinatorics \textbf{11} (2004), Dynamic Survey 11.

\bibitem{OBryant10}
K. O'Bryant,
\emph{Thick subsets that do not contain arithmetic progressions},
arXiv:0912.1494.

\bibitem{Semchankau20}
A. Semchankau,
\emph{Maximal subsets free of arithmetic progressions in arbitrary sets},
arXiv:2010.04490.

\end{thebibliography}
